\documentclass[11pt]{article}

\usepackage[margin=1in]{geometry}
\usepackage{amsmath,amssymb,amsthm}
\usepackage[hidelinks]{hyperref}

\setlength{\parindent}{0pt}
\setlength{\parskip}{0.7em}

\newtheorem*{theorem}{Theorem}
\newtheorem*{remark}{Remark}

\DeclareMathOperator{\Range}{Range}
\DeclareMathOperator{\Bdd}{B}
\newcommand{\cl}{\overline}

\begin{document}

\begin{center}
{\Large Partial Solution Packet}\\[0.25em]
{\large Nondense-square Banach algebras have a representation with non-complemented essential space}\\[0.35em]
\textbf{Source:} arXiv:1703.00882, Question 3.4
\end{center}

\section*{Source Question}

Gardella--Thiel ask whether every Banach algebra \(A\) without a left unit has
a representation \(\varphi:A\to \Bdd(X)\) on some Banach space \(X\) whose
essential space is not complemented.

The essential space of \(\varphi\) is
\[
X_\varphi=\cl{\operatorname{span}}\{\varphi(a)x:a\in A,\ x\in X\}.
\]

\section*{Result}

\begin{theorem}
Let \(A\) be a Banach algebra such that \(\cl{A^2}\ne A\). Then there is a
Banach space \(X\) and a bounded representation \(\varphi:A\to\Bdd(X)\) such
that \(X_\varphi\) is not complemented in \(X\).
\end{theorem}

\begin{proof}
Since \(\cl{A^2}\ne A\), Hahn--Banach gives a nonzero functional
\(\chi\in A^*\) such that \(\chi(ab)=0\) for all \(a,b\in A\).

Let
\[
0\longrightarrow H \stackrel{i}{\longrightarrow} Z_2
\stackrel{q}{\longrightarrow} H\longrightarrow 0
\]
be the Kalton--Peck nonsplit twisted Hilbert exact sequence
\cite{KaltonPeck}. Thus \(H\) is a Hilbert space, \(i(H)\) is a closed subspace
of \(Z_2\), \(q i=0\), and \(i(H)\) is not complemented in \(Z_2\). Put
\(X=Z_2\) and
\[
T=i q\in \Bdd(X).
\]
Then \(T^2=i q i q=0\), and \(\Range(T)=i(H)\). The range is not complemented:
if it were, the quotient map \(q:Z_2\to H\) would admit a bounded right inverse
by restricting \(q\) to a complement of \(i(H)\), contradicting nonsplitting.

Define
\[
\varphi(a)=\chi(a)T\qquad(a\in A).
\]
This is bounded and linear. Moreover, for \(a,b\in A\),
\[
\varphi(a)\varphi(b)=\chi(a)\chi(b)T^2=0
\]
while
\[
\varphi(ab)=\chi(ab)T=0.
\]
Hence \(\varphi\) is a representation.

Choose \(a_0\in A\) with \(\chi(a_0)\ne0\). Then
\[
X_\varphi
=\cl{\operatorname{span}}\{\chi(a)T x:a\in A,\ x\in X\}
=\Range(T)
=i(H),
\]
which is not complemented in \(X\). This proves the claim.
\end{proof}

\section*{Scope}

The hypothesis \(\cl{A^2}\ne A\) automatically excludes left units, since a
left unit would imply \(A=A^2\). The theorem therefore gives a positive answer
to Question 3.4 for all Banach algebras with nondense square, including all
nonzero Banach algebras with zero product.

The dense-square case remains outside this packet.

\begin{thebibliography}{9}

\bibitem{KaltonPeck}
N. J. Kalton and N. T. Peck,
\emph{Twisted sums of sequence spaces and the three space problem},
Trans. Amer. Math. Soc. 255 (1979), 1--30.

\end{thebibliography}

\end{document}
